% Generated draft: requires mathematical and source review.
\documentclass[11pt]{article}
\usepackage[margin=25mm]{geometry}
\usepackage{fontspec}
\setmainfont{Latin Modern Roman}
\usepackage{amsmath,amssymb,mathtools,unicode-math}
\setmathfont{Latin Modern Math}
\usepackage{enumitem,needspace,xurl,hyperref}
\setlength{\parindent}{0pt}
\setlength{\parskip}{5pt}
\setlength{\emergencystretch}{3em}
\newcommand{\N}{\mathbb{N}}
\newcommand{\Z}{\mathbb{Z}}
\newcommand{\Q}{\mathbb{Q}}
\newcommand{\R}{\mathbb{R}}
\newcommand{\C}{\mathbb{C}}
\newcommand{\F}{\mathbb{F}}
\newcommand{\A}{\mathbb{A}}
\newcommand{\PP}{\mathbb{P}}
\newcommand{\E}{\mathbb{E}}
\newcommand{\eps}{\varepsilon}
\newcommand{\dd}{\,\mathrm{d}}
\begin{document}
\section*{UnsolvedMath ID: 30003101}
\textbf{Source problem code:} catalogue entry
\subsection{Definitions and mathematical statement}
Let \(X\) be the smooth projective curve and \(G\) the reductive group intended in the source, and write \(\operatorname{Bun}_G(X)\) for the moduli stack of principal \(G\)-bundles on \(X\). Write \(\check G\) for the Langlands dual group, whose root datum is dual to that of \(G\) in the connected reductive case. The supplied literature assessment places the question in the finite-field \(\ell\)-adic setting: the base field is a finite field \(k\), and \(\ell\ne\operatorname{char}(k)\). The record does not specify the remaining hypotheses on \(X\) and \(G\), the coefficient field, or the descent conventions; these remain qualifications on the formulation below.

A \(\check G\)-local system \(E\) associates to a finite-dimensional representation \(V\) of \(\check G\) a local system \(E_V\) on \(X\), compatibly with tensor operations. Let \(\operatorname{Loc}_{\check G}^{\mathrm{irr}}(X)\) denote the local systems satisfying the source's intended irreducibility condition. That condition is not defined in the supplied record. In particular, whether it means that monodromy is not contained in a proper parabolic subgroup, and whether this is tested arithmetically or geometrically, must be clarified.

A Hecke eigensheaf with eigenvalue \(E\) is a sheaf complex \(F\) on \(\operatorname{Bun}_G(X)\), together with Hecke eigenisomorphisms, schematically
\[
\mathsf H_V(F)\simeq E_V\boxtimes F
\qquad\text{on }X\times\operatorname{Bun}_G(X),
\]
for every finite-dimensional representation \(V\) of \(\check G\), satisfying the tensor, unit, and iterated-Hecke compatibilities. Here \(\mathsf H_V\) denotes the Hecke functor attached to \(V\). This display presupposes compatible choices of Hecke normalization; the shifts, twists, and ambient sheaf category are not specified by the record.

\textbf{Source question, subject to these unresolved qualifications.} For every
\(E\in\operatorname{Loc}_{\check G}^{\mathrm{irr}}(X)\), does there exist a nonzero Hecke eigensheaf \(F\) on \(\operatorname{Bun}_G(X)\), with eigenvalue \(E\), such that
\[
\left.F\right|_C\text{ is a nonzero simple object of }\operatorname{Perv}(C)
\qquad\text{for every irreducible component }C\subseteq\operatorname{Bun}_G(X)?
\]
Here \(\operatorname{Perv}(C)\) denotes the intended category of perverse \(\ell\)-adic sheaves on \(C\); ``irreducible perverse sheaf'' means a nonzero simple object. The restriction requirement applies to every irreducible component, including components other than any chosen neutral component. A fully fixed mathematical statement requires the omitted hypotheses and categorical conventions to be recovered from the source.
\subsection{Short English statement}
The question asks whether every irreducible Langlands-dual local system on the curve has a nonzero Hecke eigensheaf whose restriction to every irreducible component of the bundle stack is an irreducible perverse sheaf. The supplied curation places this in the finite-field \(\ell\)-adic setting, but leaves several hypotheses and conventions unspecified.
\subsection{Sources}
\begin{itemize}
\item Supplied local UnsolvedMath registry record, permanent ID \(30003101\), title \emph{Existence of Irreducible Hecke Eigensheaves}, exact source problem code \texttt{catalogue entry}, in the Oberwolfach Reports Open Problems collection. The proposed year is recorded as 2016; no individual proposer is supplied.
\item Supplied primary-source attribution: \emph{Arbeitsgemeinschaft: The Geometric Langlands Conjecture} (2017), Oberwolfach Reports, pp.~1027--1098, DOI 10.4171/owr/2016/20. \url{https://doi.org/10.4171/owr/2016/20}. This is also the supplied literature reference labeled ``OWR 2016 report.'' The record says these reports document workshop and open-problem discussions.
\item Supplied literature reference labeled ``Gaitsgory--Raskin 2024'': \url{https://arxiv.org/abs/2409.09856}. No more detailed bibliographic attribution is supplied.
\item Literature assessment dated exactly \texttt{2026-08-21}: the supplied curation labels the problem open and states that the general reductive-group eigensheaf assertion remains unresolved in the finite-field \(\ell\)-adic setting, and that later characteristic-zero geometric Langlands results do not by themselves settle this exact claim. This is supplied dated curation, not an independently verified literature assessment.
\item The registry labels the statement \texttt{corrected\_verified} and reports that the official report passage was checked and rewritten as a self-contained problem preserving its mathematical content. That verification claim is supplied provenance; it has not been independently checked for this draft.
\end{itemize}
\end{document}
